%
%  THESISBOEK
%
%  Dit bestand zorgt voor algemene (layout)definities, en groepeert de
%  afzonderlijke LaTeX-files tot een geheel.
%
%  @author Erwin Six, David De Reu, Brecht Vermeulen
%

\documentclass[11pt,a4paper,oneside,notitlepage]{book}
\usepackage[english,dutch]{babel}
\usepackage{algorithmic}
\usepackage{algorithm}
\usepackage{amsthm}
\usepackage{hyperref}

% marges aanpassen
% (opmerking: moet *voor* inclusie van fancyhdr package komen)
\setlength{\hoffset}{-1in}
\setlength{\voffset}{-1in}
\setlength{\topmargin}{2cm}
\setlength{\headheight}{0.5cm}
\setlength{\headsep}{1cm}
\setlength{\oddsidemargin}{3.5cm}
\setlength{\evensidemargin}{3.5cm}
\setlength{\textwidth}{16cm}
\setlength{\textheight}{23.3cm}
\setlength{\footskip}{1.5cm}

\usepackage{fancyhdr}
\usepackage{graphicx}
% \usepackage[colorlinks]{hyperref}

\newtheorem{mydef}{Definition}

\pagestyle{fancy}

\renewcommand{\chaptermark}[1]{\markright{\MakeUppercase{#1}}}
\renewcommand{\sectionmark}[1]{\markright{\thesection~#1}}

\newcommand{\headerfmt}[1]{\textsl{\textsf{#1}}}
\newcommand{\headerfmtpage}[1]{\textsf{#1}}

\fancyhf{}
\fancyhead[LE,RO]{\headerfmtpage{\thepage}}
\fancyhead[LO]{\headerfmt{\rightmark}}
\fancyhead[RE]{\headerfmt{\leftmark}}
\renewcommand{\headrulewidth}{0.5pt}
\renewcommand{\footrulewidth}{0pt}

\fancypagestyle{plain}{ % eerste bladzijde van een hoofdstuk
  \fancyhf{}
  \fancyhead[LE,RO]{\headerfmtpage{\thepage}}
  \fancyhead[LO]{\headerfmt{\rightmark}}
  \fancyhead[RE]{\headerfmt{\leftmark}}
  \renewcommand{\headrulewidth}{0.5pt}
  \renewcommand{\footrulewidth}{0pt}
}

% anderhalve interlinie (opm: titelblad gaat uit van 1.5)
\renewcommand{\baselinestretch}{1.5}

% indien LaTeX niet goed splitst, neem je woord hierin op, of evt om splitsen 
% te voorkomen
\hyphenation{ditmagnooitgesplitstworden dit-woord-splitst-hier}

\begin{document}

% titelblad (voor kaft)
%  Titelblad

% Opmerking: gaat uit van een \baselinestretch waarde van 1.5 (die moet
% ingesteld worden voor het begin van de document environment)

\begin{titlepage}

\setlength{\hoffset}{-1in}
\setlength{\voffset}{-1in}
\setlength{\topmargin}{1.5cm}
\setlength{\headheight}{0.5cm}
\setlength{\headsep}{1cm}
\setlength{\oddsidemargin}{3cm}
\setlength{\evensidemargin}{3cm}
\setlength{\footskip}{1.5cm}
\enlargethispage{1cm}
% \textwidth en \textheight hier aanpassen blijkt niet te werken

\fontsize{12pt}{14pt}
\selectfont

\begin{center}

\includegraphics[height=2cm]{fig/ruglogo}

\vspace{0.5cm}

Faculteit Toegepaste Wetenschappen\\
Vakgroep Informatietechnologie\\
Voorzitter: Prof.~Dr.~Ir.~Dani\"{e}l~de~Zutter

\vspace{3.5cm}

\fontseries{bx}
\fontsize{17.28pt}{21pt}
\selectfont

Identifying experts through a framework for\\
knowledge extraction from public online source

\fontseries{m}
\fontsize{12pt}{14pt}
\selectfont

\vspace{.6cm}

by 

\vspace{.4cm}

Simon Buelens, Mattias Putman

\vspace{3.5cm}

Promotors: Prof.~Dr.~Ir.~Filip~De~Turck,~Elena~Tsiporkova~(Sirris)\\
Scriptiebegeleiders: Anna~Hristoskova,~Tom~Tourw\'{e}~(Sirris)\\

\vspace{2cm}

Masterproef ingediend tot het behalen van de academische graad van\\
Master in de ingenieurswetenschappen: computerwetenschappen

\vspace{1cm}

Academiejaar 2010--2011

\end{center}
\end{titlepage}


% lege pagina (!!)

% titelblad (!!)

% geen paginanummering tot we aan de inhoudsopgave komen
\pagestyle{empty}

% voorwoord met dankwoord en toelating tot bruikleen (ondertekend)
%  Voorwoord (dankwoord) en toelating tot bruikleen

\newpage

\noindent \textbf{\huge Dankwoord}

\vspace{1.5cm}

\noindent
We willen vooreerst de gelegenheid aangrijpen om een aantal mensen te bedanken. 
\\

\noindent
We danken onze begeleiders Elena Tsiporkova, Tom Tourw\'e, Anna Hristoskova en Tim Wauters voor de goeie inzichten, de aangename vergaderingen en hun eeuwige geduld.
\\

\noindent
Daarnaast willen we Anne-Sophie Delbecque immens veel bedanken voor het annoteren van de testset, een tijdrovend en eentonig werk.
\\

\noindent
Tot slot bedanken we ook onze families voor de goede zorgen en het vele geduld gedurende de volledige studies in het algemeen en deze moeilijke periode in het bijzonder.

\vspace{1.5cm}

\noindent

\addvspace{4cm}

\noindent Simon Buelens en Mattias Putman, augustus 2011\newpage

\noindent \textbf{\huge Toelating tot bruikleen}

\vspace{1.5cm}

\noindent
``De auteur geeft de toelating deze scriptie voor consultatie beschikbaar
te stellen en delen van de scriptie te kopi\"eren voor persoonlijk
gebruik.\\
Elk ander gebruik valt onder de beperkingen van het auteursrecht,
in het bijzonder met betrekking tot de verplichting de bron uitdrukkelijk
te vermelden bij het aanhalen van resultaten uit deze scriptie.''

\addvspace{4cm}

\noindent Simon Buelens en Mattias Putman, augustus 2011\newpage


% overzicht
%  Overzichtsbladzijde met samenvatting

\newpage

{
\setlength{\baselineskip}{32pt}
\setlength{\parindent}{0pt}
\setlength{\parskip}{18pt}

\begin{center}

\noindent \textbf{\huge
Identifying experts through }
\textbf{\huge a framework for knowledge extraction}
\textbf{\huge from public online sources}

\setlength{\baselineskip}{12pt}
\setlength{\parindent}{0pt}
\setlength{\parskip}{12pt}

door 

Simon Buelens, Mattias Putman

Promotors: Prof.~Dr.~Ir.~Filip~De~Turck,~Elena~Tsiporkova~(Sirris),~Tom~Tourw\'{e}~(Sirris)\\
Scriptiebegeleiders: Anna~Hristoskova,~Tim~Wauters\\

Masterproef ingediend tot het behalen van de academische graad van\\
Master in de ingenieurswetenschappen: computerwetenschappen

Academiejaar 2010-2011\\
Faculteit Ingenieurswetenschappen\\
Voorzitter: Prof. Dr. Ir. Dani\"{e}l De Zutter\\
Vakgroep Informatietechnologie\\

\end{center}

\setlength{\baselineskip}{10pt}
\setlength{\parindent}{0pt}
\setlength{\parskip}{10pt}

\renewcommand{\baselinestretch}{1.1} 	% De interlinie afstand wat vergroten.
\small\normalsize                       % Nodig om de baselinestretch goed te krijgen.

\section*{Samenvatting}

Onderzoekers verliezen veel tijd met de zoektocht naar informatie gerelateerd aan hun onderzoeksdomein. Er bestaan bijna geen diensten die toelaten om aan de hand van trefwoorden een overzicht te verkrijgen met experts voor de opgegeven domeinen. Er is onderzoek naar disambiguatie van auteurs, maar deze worden meestal niet in combinatie gebracht met het opzoeken van experten, maar het indelen van publicaties (alhoewel de twee gerelateerd zijn).

In deze masterproef onderzoeken we de mogelijkheid om een framework op te stellen dat dit toelaat door online informatie op te zoeken, deze informatie in relatie te brengen met de correcte auteurs en gebruikers toe te laten dit framework te gebruiken om hierin te zoeken. De nadruk van het framework ligt op de disambiguatie van auteurs (aan de hand van de aanwezige informatie alle namen zo goed mogelijk connecteren met de juist auteur) aan de hand van een regelgebaseerde aanpak en de uitbreidbaarheid van het framework.

We maken gebruik van een graafgebaseerde representatie van de data en de architectuur is gebaseerd op pipes en filters. Dit laat toe dat het framework uitbreidbaar, schaalbaar en eenvoudig aanpasbaar is. Op het einde volgen de resultaten gebaseerd op een manueel geannoteerde testset. Uiteindelijk gaan we ook de vergelijking aan met de verdeling van auteurs door DBLP.

\section*{Trefwoorden}

auteur disambiguatie, data verwerking, clustering, pipes en filters

}

\newpage % strikt noodzakelijk om een header op deze blz. te vermijden


\pagestyle{fancy}
\frontmatter

% inhoudstafel
\tableofcontents

% opmaak voor het eigenlijke boek; onderstaande lijnen
% weglaten als de eerste regel van een nieuwe alinea moet
% inspringen in plaats van extra tussenruimte
%\setlength{\parindent}{0pt}
%\setlength{\parskip}{0.5\baselineskip plus 0.5ex minus 0.2ex}
%\setlength{\parskip}{1ex plus 0.5ex minus 0.2ex}

% hoofdstukken
\mainmatter

% hier worden de hoofdstukken ingevoegd (\includes)
\include{hoofdstukx}

% appendices
\appendix

% hier worden de appendices ingevoegd (\includes)


\include{referenties}

\backmatter

% eventueel: lijst van figuren en tabellen
%\listoffigures
%\listoftables

% lege pagina (!!)

% kaft

\end{document}
